% CONTRACT
%   Every subsection states which table it reads. A number with no table is a
%   defect, not a draft. Resolve via paper/bmc/_tables.py (newest_by_time).
%   Quote weighting-dependent statistics WITH their weighting and denominator.

\subsection{Interval methods are miscalibrated under the tested generators}
\label{sec:interval-stress}
% TABLE: interval_stress_calibration_summary (126), _calibration (630)
% Contract kept: a result about these TESTED GENERATORS, never a general claim.
% Numbers re-derived in tests/test_bmc_paper_numbers.py.

We varied only the distribution of the per-patient effect, holding the Poisson
thinning, the cohorts and the nominal level fixed. Seven cohorts, six generators,
three interval methods, five fixed seeds and 800 trials per cell give 630 cells,
each with a binomial standard error and a Wilson interval. Every regime is a true
null --- the effect has mean zero --- so a rejection is a false positive and the
measured rate is a false-positive rate.

The gaussian regime is the generator the committed calibration used, and it acts
as the control: it reproduces the committed rates. Percentile over-rejection is
$9.9\%$ on the MLH1-methylated arm against the committed $10.7\%$, and $6.1\%$ on
the adenoma lineage rung against $5.7\%$.

Under every non-saturated generator the percentile bootstrap and BCa over-reject
(Table~\ref{tab:stress}). The Student-$t$ interval is calibrated in the same
cells, with one exception: under skewed effects on the MLH1-methylated arm
($n=10$), where its median false-positive rate reaches $7.0\%$. The boundary
regimes separate two mechanisms. Where detection saturates every method is
conservative ($0.0\%$), not because the interval is good but because both arms
detect every cell, so the per-patient deltas are identically zero; ordinary input
validation fires on $35$ of $35$ such cells. In the interior it is nearly blind,
firing on at most one of $35$ cells for any method under gaussian, skewed,
heavy-tailed or spike-slab effects. Where detection is near zero it fires on at
most $3$ of $35$ cells, while the audit fires on $17$ of $35$ percentile cells.

\begin{table}[h]
\centering
\caption{Median false-positive rate (\%), by interval method and generator,
across seven cohorts, five seeds and 800 trials per cell. Nominal is $5\%$. Every
regime is a true null. The gaussian column is the committed generator and the
control.}
\label{tab:stress}
\begin{tabular}{lrrrrrr}
\toprule
 & gaussian & skewed & heavy-tailed & spike-slab & boundary-rare & boundary-sat. \\
\midrule
percentile  & 7.0 & 7.9 & 7.8 & 6.6 & 6.4 & 0.0 \\
BCa         & 7.1 & 7.8 & 8.6 & 7.1 & 7.1 & 0.0 \\
Student-$t$ & 4.8 & 5.9 & 4.9 & 4.4 & 4.9 & 0.0 \\
\bottomrule
\end{tabular}
\end{table}

The pre-audit diagnostic --- compare the percentile rate to the closed-form
normal approximation --- catches skew ($12$ of $35$ percentile cells) but misses
heavy tails ($3$ of $35$), because the closed form is a function of $n$ alone and
carries no information about the effect distribution. Across the seven cohorts
the measured percentile rate tracks the closed form closely under gaussian
effects ($19.1\%$ against $18.8\%$ at $n=4$) and departs from it under skew.

\textbf{What this does and does not establish.} This is a result about the tested
generators (decision record \S3). It does not establish that BCa is generally
worse, that the width expression is an exact coverage formula, or that
Student-$t$ is universally calibrated.

\subsection{Reported cutpoints are brackets, and one rule silently drops a
crossing}
\label{sec:brackets}
% TABLE: cutpoint_crossing_brackets_summary (8), _brackets (64)
% Contract kept: report a bracket, never a precise threshold.

A calibration returns the smallest bin-median mature-cell count that meets a
criterion. That is a point on a coarse grid, and the crossing it reports lies
somewhere between the last failing and the first passing bin. We report that
interval instead of the point. The two criteria are kept apart: \emph{wide} is
coverage alone; \emph{ok} is coverage and discrimination.

The separation is what exposes a silent loss. The point rule requires
\emph{ok}, so when no bin is \emph{ok} it writes \emph{wide} as None. On the
pooled draw pool that is exactly what happens, and a coverage crossing exists
anyway: on both cohorts the \emph{wide} bracket is $[42.5, 70]$ mature cells,
stable across all eight seeds, while \emph{ok} is \notest{} on all eight. The
bracket table recovers a crossing the point rule discards.

The cohorts are reported separately and never pooled. SMC's reference pool gives
\emph{wide} $[27.5, 42.5]$ and \emph{ok} $[42.5, 70]$, both stable. KUL3's
reference \emph{wide} is $[42.5, 70]$, stable, while its \emph{ok} bracket is
unstable: $[150, 400]$ on one seed and $[400, 800]$ on the others. A precise
threshold would hide that spread; the bracket is the reported quantity.

\subsection{The descriptive adenoma pattern, and what fixing each choice costs}
\label{sec:adenoma}
% TABLE: adenoma_claim_sensitivity_summary (90), _sensitivity (2700),
%        estimability_attrition (144)
% Contract kept: no silencing claim; statistics quoted with weighting.

Recomputing the adenoma summaries across all three weightings, four rungs, both
denominators and all four scale-free statistics gives 2700 contrast rows and 144
attrition rows. The primary (resolved) portion is bit-identical to the committed
scale-free table: all 1260 rows agree on centre, interval, verdict and patient
count.

The denominator is not cosmetic. Under the primary (resolved) denominator the
load-bearing statistic, $\log(|i|/|c|)$, is undefined on 792 of the epithelial
rows, because the compositional term is exactly zero there by construction; under
the rejected all-epithelial denominator only 4 rows are undefined. At the lineage
rung, where the descriptive claim is made, the load-bearing statistic separates
all 8 cross-block contrasts from zero under every weighting and both
denominators. Requiring all four statistics to agree instead, 6 of 8 cross-block
contrasts survive under doubly-robust and normal weighting and 5 of 8 under
tumour weighting. \textbf{Those are the numbers to quote, with their weighting;
on the load-bearing statistic the two target genes are not distinguishable at any
rung, and that lack of a distinction does not establish their equivalence}, so no
gene-specific claim follows.

Attrition is reported rather than assumed. At lineage the resolved denominator
retains 44 patients, one of whom is \notest{}; at best4 it retains 20, two of
whom are \notest{}. We do not report a share of loss as a share of variance: the
mean absolute intrinsic share is $|i|/(|i|+|c|)$, which discards signs and
excludes the interaction from its denominator.

\subsection{Two reproductions from committed derived inputs}
\label{sec:reproductions}
% TABLE: crowell_summary_reproduction (5), disval_reproduction (64)
% Contract kept: labelled reproduction from derived inputs, not replication.

Two readings whose raw material is no longer on disk are reproduced from derived
tables already in version control, and labelled as such. The Crowell spatial
summary is recomputed from the committed per-section feasibility tables: all 5
genes reproduce exactly against the canonical run, with a zero difference in mean,
interval, block count and interval verdict. The DIS/VAL stability split is
rebuilt from the committed patient-to-collection assignment rather than the
interim atlas cache, and all 64 cells (4 halves $\times$ 4 statistics $\times$ 4
contrasts) reproduce exactly.

Neither is a raw-data replication. The Crowell result remains supporting
evidence for a tissue-level pattern and is not an independent replication of the
decomposition or of its mechanism.

\subsection{A resolver that cannot say which run it read}
\label{sec:resolver}
% TABLE: table_resolution_ambiguity (23 rows). This is the audit's own finding.

Versioned result directories are named \texttt{\{date\}\_\{sha\}}, and the
project's resolver took the lexicographically last, which orders runs within a
date by commit hash rather than by time. Enumerating every table with more than
one same-date run and comparing the two candidate frames: 23 table/date pairs
disagree between name order and time order, and 9 of them hold different content.
It is not hypothetical --- it produced a wrong Crowell comparison in this audit
before it was caught. The affected \texttt{icbi\_coexpression} table feeds a
calibration presented above, but its Pelka rows are byte-identical between the
candidates, so no published number is wrong.

Neither rule is universally right. For the panel-coverage table and the
lesion-subtype falsifier branch the name-order pick is canonical; for Crowell,
Becker, the MLH1 power tables and the interval calibration the time-order pick is.
The repair is therefore not a blanket resolver change but an explicit one:
resolve by sidecar timestamp where a comparison needs the latest run, and record
the ambiguity as a table. This finding came from the audit running on the audit's
own infrastructure, which is the strongest available evidence that the failure
mode is structural rather than careless.

\subsection{What the audit could not evaluate}
\label{sec:not-evaluated}
% TABLE: blinded_guard_challenge (12 rows), the blocked rows of
%        reproduction_manifest. This is why the audit is not yet validated.

The blinded challenge has not been executed. The committed set of 12 cases (6
clean controls) was constructed by a worker who had read the ledger, so every
case carries \texttt{saw\_ledger=true}; \texttt{is\_held\_out\_evaluation} is
false, and the week-one stop rule is not readable from the table. In the
reproduction index both the challenge and its results table carry status
\textsc{blocked} rather than \textsc{present}. The blinded challenger and the
second ledger rater are unfilled.

Until that is discharged, the audit's stop rule --- withdraw the methods
submission if no additional decision-relevant failure is found --- cannot be
applied, and this manuscript cannot be submitted. This is stated in the main text
rather than buried in the Discussion because it is the reason the reader cannot
yet take the audit as validated.
